% Generated by careport.psg_assist_manuscript from supplement.aux; do not edit.
\makeatletter
\newlabel{tab:s-codes}{{S1}{6}{Names used in the text and the codes used in the analysis plans (Supplementary File S1)}{table.caption.11}{supplement.pdf}}
\newlabel{tab:s-r1r2}{{S2}{6}{Acceptance lift by the pre-fill's sleep stage (R1) and anchoring by respiratory event type (R2)}{table.caption.12}{supplement.pdf}}
\newlabel{tab:s-r3r4}{{S3}{6}{Dependence on ambiguity (R3) and adoption of pre-fill labels that most experts rejected manually: selected with the scored expert's own label (R4) and without it (C1)}{table.caption.13}{supplement.pdf}}
\newlabel{tab:s-r6}{{S4}{7}{Robustness of the primary conclusions to analysis choices (R6, C6)}{table.caption.14}{supplement.pdf}}
\newlabel{tab:validation}{{S5}{7}{Validation of the two model configurations against observations they were not fitted to (sleep stages, percent agreement, pooled over epochs)}{table.caption.15}{supplement.pdf}}
\newlabel{tab:consequences}{{S6}{7}{Consequences in the calibrated model: references built from manual versus assisted scorings}{table.caption.16}{supplement.pdf}}
\newlabel{tab:s-sensitivity}{{S7}{7}{Sensitivity analyses of the simulation: 12 assisted experts, majority consensus}{table.caption.17}{supplement.pdf}}
\newlabel{tab:s-ranking}{{S8}{8}{Probability that a truly better competitor is ranked above the pre-fill (calibrated model)}{table.caption.18}{supplement.pdf}}
\newlabel{tab:s-fda}{{S9}{8}{How the 22 FDA clearances of automatic PSG scoring software (product code OLZ, 2012 to 2026) with clinical performance data built their reference, as stated in the public 510(k) summaries}{table.caption.19}{supplement.pdf}}
\newlabel{tab:s-predictions}{{S10}{9}{Pre-specified predictions and outcomes}{table.caption.20}{supplement.pdf}}
\newlabel{tab:s-stagers}{{S11}{10}{Negative control and comparisons with the two stagers that no expert saw (sleep stages; C2, C3, C5)}{table.caption.21}{supplement.pdf}}
\newlabel{tab:s-indices}{{S12}{10}{Clinical indices (events per hour of sleep) of the pre-fill and the two consensuses (M4)}{table.caption.22}{supplement.pdf}}
\newlabel{tab:s-m3}{{S13}{10}{Agreement among the 12 experts where the pre-fill agreed or disagreed with the manual consensus (M3)}{table.caption.23}{supplement.pdf}}
\newlabel{tab:s-blocks}{{S14}{11}{Audits of 30-minute blocks at five sizes, and audits of single epochs (sleep stages, C4 and B1)}{table.caption.24}{supplement.pdf}}
\newlabel{tab:s-units}{{S15}{11}{The 20 recording-task units: anchoring, the two consensus changes, the difference-in-differences and side-taking (M1, M2, A2)}{table.caption.25}{supplement.pdf}}
\newlabel{tab:s-experts}{{S16}{12}{Anchoring by expert: each expert's mean $\Delta A$ over the five recordings of each task (M1)}{table.caption.26}{supplement.pdf}}
\newlabel{tab:s-d1}{{S17}{12}{Blind second coding of the 22 FDA summaries with performance data (amendment 005, D1): agreement of each coder with the first coding, overall and by the coder's exposure}{table.caption.27}{supplement.pdf}}
\newlabel{tab:s-d1-items}{{S18}{12}{The 8 fields on which a coder's code differed from the first coding, in 5 summaries, and their adjudication from the summary text (amendment 005, D1)}{table.caption.28}{supplement.pdf}}
\newlabel{fig:s-simulation}{{S1}{13}{Simulation with the truth known (calibrated model; outputs are illustrations under the model's conditions). (a) The pre-fill's measured minus true accuracy against references of 1 to 25 experts (majority consensus), for manual and assisted references. (b) The handicap in favor of the pre-fill over an equally accurate competitor, against anchoring strength, for references of 1, 12 and 25 assisted experts; the dashed line marks the calibrated $\alpha $. (c) Model against observed values for the validation targets}{figure.caption.30}{supplement.pdf}}
\makeatother
